\section{Постановка задачи о потенциальном обтекании сферы}

\subsection{Геометрия задачи}
Рассмотрим задачу о потенциальном обтекании сферы радиусом \( R \), находящейся в потоке идеальной несжимаемой жидкости. Поток набегает на сферу вдоль оси \( z \) с постоянной скоростью \( U_\infty \). В силу осевой симметрии задачи удобно использовать цилиндрическую систему координат \( (r, \phi, z) \), где ось \( z \) совпадает с направлением набегающего потока, а \( r \) и \( \phi \) — радиальная и азимутальная координаты соответственно.

\subsection{Осесимметричное приближение}
В силу осевой симметрии задачи все величины не зависят от азимутальной координаты \( \phi \). Таким образом, задача сводится к двумерной в плоскости \( (r, z) \).

\subsection{Уравнение для потенциала скорости}
Потенциал скорости \( \Phi(r, z) \) удовлетворяет уравнению Лапласа в цилиндрических координатах:
\[
\frac{1}{r} \frac{\partial}{\partial r} \left( r \frac{\partial \Phi}{\partial r} \right) + \frac{\partial^2 \Phi}{\partial z^2} = 0.
\]
Это уравнение следует из условия несжимаемости \( \nabla \cdot \mathbf{v} = 0 \) и потенциальности течения \( \mathbf{v} = \nabla \Phi \).

\subsection{Функция тока}
Функция тока \( \Psi(r, z) \) вводится для удобства описания течения и связана с компонентами скорости соотношениями:
\[
v_r = \frac{1}{r} \frac{\partial \Psi}{\partial z}, \quad v_z = -\frac{1}{r} \frac{\partial \Psi}{\partial r}.
\]
Функция тока также удовлетворяет уравнению Лапласа:
\[
\frac{\partial^2 \Psi}{\partial r^2} - \frac{1}{r} \frac{\partial \Psi}{\partial r} + \frac{\partial^2 \Psi}{\partial z^2} = 0.
\]

\subsection{Граничные условия}
На поверхности сферы выполняется условие непротекания:
\[
\left. \frac{\partial \Phi}{\partial n} \right|_{r^2 + z^2 = R^2} = 0,
\]
где \( n \) — нормаль к поверхности сферы. На бесконечности поток стремится к однородному:
\[
\nabla \Phi \to U_\infty \mathbf{e}_z \quad \text{при} \quad r^2 + z^2 \to \infty.
\]
Для функции тока граничные условия формулируются аналогично.